%% ma: L10-B14-0082
%% lop: 10
%% chuong: 5
%% bai: 14
%% dang: 10.14.8
%% loai: TN4
%% muc: VD
%% dap_an: "C"
%% nguon: "chuyen-de-cac-so-dac-trung-cua-mau-so-lieu-khong-ghep-nhom-toan-10-chuong-trinh-moi.pdf (GV. Phan Nhật Linh) (P05), Bài 03 Dạng 2 Phần I Câu 29"
%% kien_thuc: "Lớp 10 Bài 13 (số trung bình), Bài 14 (phương sai, độ lệch chuẩn)"
%% trang_thai: chinh-thuc
%% ghi_chu: "Tự duyệt theo yêu cầu giáo viên 22/10/2026."
\begin{ex}
Kết quả thi học kì I môn Toán của $48$ học sinh lớp 10A được cho bởi bảng tần số sau:
\begin{center}
\begin{tabular}{|c|c|c|c|c|c|c|}\hline
Điểm & $5$ & $7$ & $8$ & $8{,}5$ & $9$ & $10$\\\hline
Tần số & $1$ & $3$ & $12$ & $4$ & $20$ & $8$\\\hline
\end{tabular}
\end{center}
Phương sai và độ lệch chuẩn của mẫu số liệu lần lượt gần với kết quả nào nhất?
\choice{$8{,}67$ và $0{,}91$}{$0{,}91$ và $0{,}83$}{\True $0{,}91$ và $0{,}95$}{$0{,}91$ và $0{,}46$}
\loigiai{$\overline x=\dfrac{416}{48}\approx8{,}67$; $s^2=\dfrac1{48}\sum n_i\left(x_i-\overline x\right)^2\approx0{,}91$; $s\approx0{,}95$. Chọn C.}
\end{ex}
